\chapter{Covalent Bonds: Unequal Electron Sharing}

\begin{kaichang}
Hospitals have blue steel cylinders containing compressed oxygen. Mountaineers carry them up Everest, while divers carry cylinders containing oxygen mixed with other gases. About a fifth of the breath you are taking now is oxygen, floating around as ``two oxygen atoms joined together'' and passing from your alveoli into your blood.

Now look at some other things nearby: a plastic pen barrel, bottled water, a rubber band, your own hair. They differ from oxygen and from salt. Salt shatters when dropped, disappears in water, and conducts when molten (Chapter 18). Plastic deforms before breaking when pulled; water boiled into vapor is still water; burning hair smells charred instead of simply ``melting.''

Now guess: why must an oxygen atom stick to another oxygen atom? Two identical atoms have identical electron-grabbing strength, so salt's ``one gives, the other takes'' approach (Chapter 18) cannot work. How do they join? By the end of this chapter, we will use a new model to explain that blue cylinder, that glass of water, and that plastic ruler anew.
\end{kaichang}

\section{When ``Giving Up an Electron'' Is Not an Option}

Recall Chapter 18's ionic bonding. Sodium has only 1 outer electron and loses it ``cheaply''; chlorine has 7 and fills its outer shell by taking 1. One gives, the other grabs, leaving both outer arrangements more comfortable and lowering the whole system's energy. Electrons transfer, and Coulomb attraction packs the resulting ions into a crystal. This account balances only if \textbf{the partners' appetites for electrons differ enough}.

But many everyday substances are not built this way. Hydrogen with hydrogen, oxygen with oxygen, chlorine with chlorine: when identical atoms meet, neither can take an electron for free. Each pulls as strongly as the other, producing a stalemate. Carbon with hydrogen, carbon with oxygen, or nitrogen with hydrogen is similar. Their appetites differ, but not enough for either to confiscate the other's electron outright.

Yet these substances certainly exist stably. Two hydrogen atoms bind firmly into a hydrogen molecule, \chem{H_2}, requiring considerable energy to separate. Visible evidence is everywhere: hydrogen can sit quietly in air (provided nothing ignites it), oxygen stays in cylinders for years, and plastic bags hold heavy loads without falling apart. Clearly, atoms have another way of joining, quite unlike Chapter 18's.

Compare the observable facts and the differences are obvious. Salt melts above $800\,^{\circ}\mathrm{C}$, insulates as a solid, and conducts when molten or dissolved. Sugar, made of carbon, hydrogen, and oxygen, starts changing color and decomposing around $190\,^{\circ}\mathrm{C}$ without even giving you a chance to melt it, and its aqueous solution barely conducts. Wax softens readily; ice becomes water when heated. Their common feature is that their ``basic teams'' are separate little groups---\textbf{molecules}. Atoms within each group bind firmly, while neighboring groups attract only weakly (Chapter 22's topic). Chapter 18's ``electron transfer plus electrostatic packing'' model cannot explain this internal firmness.

\section{Sharing: Two Nuclei Hold the Same Electron Pair}

Zoom in on two approaching hydrogen atoms, each with one proton and one electron. Chapter 17 set the stage: nuclei repel nuclei and electrons repel electrons, but each electron also feels attraction from the other nucleus. Total energy varies with distance along a potential-energy curve that first falls and then rises. Its ``valley bottom'' gives the bond length, and its depth the bond energy.

Now the crucial question: why does this valley exist for two hydrogen atoms? Why not repulsion all the way, with energy rising as they approach?

The answer lies in electrons' ``dual membership.'' As the atoms approach closely enough, the regions available to their electrons overlap. \textbf{Each electron can now feel both nuclei's attraction simultaneously}. In Chapter 9's electron-cloud language, the clouds overlap and the probability of finding electrons between the nuclei increases markedly. Electrons cease to be ``private property'' and become jointly ``rented'' by both nuclei. Attraction to two positive charges gives an electron lower energy than attraction to one nucleus alone. Together, these attractions outweigh nucleus--nucleus and electron--electron repulsion, creating the valley. This jointly occupied electron pair is a \textbf{covalent bond}.

Notice the contrast with Chapter 18. In ionic bonding, ownership transfers: sodium's electron moves to chlorine, and the resulting separate ions attract electrostatically. In covalent bonding, ownership is shared; both nuclei's attraction to the pair binds them together, with neither able to monopolize it. An imperfect but memorable analogy: ionic bonding resembles selling a house and taking the money; covalent bonding resembles two families buying one house together, whose arrangement falls apart if either leaves.

The same reasoning immediately explains many substances. Two chlorine atoms each contribute 1 electron, apparently completing both outer octets and forming \chem{Cl_2}. Oxygen shares a pair with each of two hydrogens to make water, \chem{H_2O}. Carbon has 4 outer electrons and can share one pair with each of four hydrogens to form methane, \chem{CH_4}, the principal component of natural gas burned in your stove. Carbon is especially versatile: it shares with other carbons to build chains and rings, creating millions of molecules. That is the story of Part Five (starting in Chapter 36), and every building block in that great drama is a covalent bond.

\section{One Pair, Two Pairs, Three Pairs}

Atoms need not share just one pair of electrons. Two atoms can share two or even three:

\begin{itemize}[itemsep=2pt]
\item \textbf{Single bond}: one shared pair. The bond between two hydrogen atoms and the O--H bonds in water are single bonds.
\item \textbf{Double bond}: two shared pairs. Oxygen's two atoms have a double bond; in exhaled carbon dioxide, \chem{CO_2}, carbon has a double bond to each oxygen.
\item \textbf{Triple bond}: three shared pairs. The most abundant molecule in air, nitrogen, \chem{N_2}, has a triple bond between its two atoms.
\end{itemize}

The number of shared pairs is also the \textbf{bond order}: 1 for a single bond, 2 for a double, 3 for a triple. Higher bond order makes the electron cloud between nuclei ``thicker,'' binding them more tightly. This is a measurable pattern, not a slogan. Here are real data:

\begin{table}[h]
\centering
\begin{tabular}{lccc}
\toprule
Bond & Order & Length (\ud{10^{-12}}{m}) & Energy (\ud{}{kJ\,mol^{-1}}) \\
\midrule
\chem{N-N} (single) & 1 & 145 & 163 \\
\chem{N=N} (double) & 2 & 125 & 418 \\
\chem{N \equiv N} (triple) & 3 & 110 & 945 \\
\chem{C-C} (single) & 1 & 154 & 347 \\
\chem{C=C} (double) & 2 & 134 & 614 \\
\chem{O=O} (in oxygen) & 2 & 121 & 498 \\
\chem{H-H} (in hydrogen) & 1 & 74 & 432 \\
\bottomrule
\end{tabular}
\caption{Lengths and energies of several covalent bonds. Bond energy is the energy absorbed to break 1 mole of that bond (breaking absorbs energy; forming releases it, as in Chapter 17). Values are approximate; the same bond type varies slightly between molecules.}
\end{table}

The trend is clear: higher bond order means shorter, stronger bonds. Nitrogen's triple bond is about a quarter shorter than its single bond, yet nearly six times as strong.

\begin{qiaoliang}
How large is bond energy? Consider hydrogen. Breaking 1 mole of H--H bonds requires \ud{432}{kJ}, and a mole contains \ud{6.02\times 10^{23}}{} molecules, so breaking \textbf{one} hydrogen molecule requires on average:
\[432\times 10^{3}\div 6.02\times 10^{23}\approx 7.2\times 10^{-19}\ \mathrm{J}\]
The average energy of molecular thermal motion at room temperature is only about $4\times 10^{-21}$ J---nearly two hundred times less. That is why molecules collide continually at room temperature without breaking covalent bonds: the money handed over in a collision falls far short of the ``breakup fee.'' Nitrogen's triple-bond energy is \ud{945}{kJ\,mol^{-1}}, or about $1.6\times 10^{-18}$ J per bond, more than twice hydrogen's fee. This makes nitrogen exceptionally ``lazy'': almost 80 percent of the air passes you daily while joining almost no reactions. Food factories exploit this by filling crisp packets with nitrogen as a protective gas: cheap, nontoxic, and reluctant to react.
\end{qiaoliang}

\section{An Unequal Tug-of-War: Shifting the Electron Cloud}

So far, examples such as hydrogen and chlorine molecules have identical partners with equal strength, leaving the shared pair unbiased in the middle. This is a \textbf{nonpolar covalent bond}.

When different atoms bond, however, perfectly fair sharing is almost impossible. Chapter 11 introduced \textbf{electronegativity}: an atom's ability to pull shared electrons toward itself when bonding. Fluorine ranks highest, followed closely by oxygen, nitrogen, and chlorine; electronegativity generally decreases toward the periodic table's lower left. When atoms with different electronegativities share a pair, its cloud shifts toward the more electronegative partner, ``spending more time'' on that side of the shared region.

What follows? The more electronegative end gains electron density and a small negative charge (written $\delta^{-}$). The other loses some density and becomes slightly positive ($\delta^{+}$). The Greek letter $\delta$ here denotes a ``small, incomplete'' charge, far less than an ion's whole unit. This unevenly distributed cloud makes a \textbf{polar covalent bond}.

Water is the best example. Oxygen's electronegativity (3.4) substantially exceeds hydrogen's (2.1), pulling the O--H electron cloud toward oxygen: oxygen carries $\delta^{-}$, hydrogen $\delta^{+}$. A water molecule has one slightly negative end and two slightly positive ones. Remember this picture: it will make waves in Chapters 22 and 23. Nearly all of water's quirks---its unusually high boiling point, floating ice, and remarkable dissolving power---begin with this ``little shift.''

There is a deeper insight: \textbf{nonpolar covalent, polar covalent, and ionic bonds are not three separate kinds of thing, but positions on one continuous slide}. With an electronegativity difference of 0 (identical atoms), sharing is perfectly fair at the slide's top. As the difference grows, the cloud shifts further, producing polarity. Beyond a certain difference (empirically about 1.7), the electron is taken outright and reaches the bottom: Chapter 18's ionic bond. Sodium and chlorine differ by about 2.2, a standard ionic bond; water's O--H difference is about 1.4, a distinctly polar covalent bond; hydrogen chloride, \chem{HCl}, differs by about 0.9, giving moderate polarity. Nature draws no line declaring ``only below here is covalent.'' Categories are convenient human labels; the slide itself is smooth.

\begin{shiyan}
\textbf{A Bending Stream of Water.} Materials: a plastic comb (or ruler), your hair, and a kitchen tap. Procedure: adjust the tap to a thin, continuous stream---as thin as possible without breaking up. Rub the comb vigorously through dry hair a dozen or more times, then slowly approach the stream without touching it. What to observe: the water bends visibly toward the comb, as though pulled by an invisible hand. Try more rubbing and combs made of different materials, comparing the deflection. Explanation: rubbing charges the comb. Water molecules have slightly positive and negative ends; near the charged comb they turn and align with their opposite-charge ends facing it, pulling the whole stream sideways. If electrons in water were truly ``shared equally,'' without polarity, the stream would ignore the comb. Safety note: keep away from outlets and appliances, and dry the worktop afterward. That is all; this experiment is very safe.
\end{shiyan}

\section{Drawing Molecular Plans: Lewis Structures}

After all that explanation, it is finally time for symbols. In 1916, the American chemist Lewis invented a drawing method still used on blackboards worldwide: show only each atom's \textbf{outermost electrons} as dots (or crosses). A pair shared between atoms was later simplified to a short line. These diagrams are \textbf{Lewis structures}.

There are only three steps. First, count all the atoms' outer electrons: the ``total budget.'' Second, connect the atoms with single bonds. Third, distribute the remaining electrons as ``lone pairs,'' aiming to give each atom eight surrounding electrons (except hydrogen, which needs only 2). If there are not enough, ``promote'' a lone pair to sharing, producing a double or triple bond.

Practice with three old friends. Hydrogen: 2 electrons total, exactly one pair, written H--H. Water: oxygen brings 6 and the two hydrogens 1 each, totaling 8. Two O--H single bonds use 4; the remaining 4 form two lone pairs on oxygen. Oxygen has 8 around it and each hydrogen 2, using the budget exactly. Carbon dioxide: carbon's 4 plus 6 from each oxygen gives 16. Two initial single bonds use 4, but however you distribute the rest, carbon cannot get its octet. Bring one lone pair from each oxygen into sharing to make two C=O double bonds, and all 16 electrons fit.

This accounting method also tracks a finer balance, \textbf{formal charge}. Compare the electrons assigned to an atom (all its lone-pair electrons and half its shared electrons) with the outer electrons it \textbf{brought itself}; the difference is its formal charge. This asks, ``Under this ownership plan, does the atom have more or fewer electrons than when alone?'' Consider unusual carbon monoxide, \chem{CO}. Its 10 outer electrons can give both atoms octets only as a triple bond plus two lone pairs. Oxygen receives 3 electrons from the triple bond and 2 from its lone pair, totaling 5---one fewer than the 6 it brought, so its formal charge is $+1$. Carbon is assigned 5, one more than its original 4, giving $-1$. Counterintuitively, \textbf{the more electronegative oxygen has a positive formal charge}. Formal charge is not the actual charge distribution (the real cloud still favors oxygen; CO's polarity is small, but oxygen remains slightly negative). It is only an accounting tool, but a useful one: reasonable Lewis structures minimize formal charges and preferably place negative ones on more electronegative atoms. This selects plausible drawings far better than guessing.

Some molecules cannot fit into a single Lewis structure. Ozone, \chem{O_3}, has 18 outer electrons and must be drawn with one double and one single bond---but which side gets the double bond? Both drawings are valid. Experimentally, its two O--O bonds are \textbf{exactly the same length}, about \ud{128}{10^{-12}m}, between single and double. The solution is that the real molecule is neither drawing but their ``average'': two extra electron pairs are \textbf{delocalized} over all three oxygens, giving each bond an order of ``one and a half.'' Two Lewis drawings placed side by side with a double-headed arrow are called \textbf{resonance structures}. Remember: resonance does not mean the molecule switches back and forth. It has one unique real structure throughout; our paper drawings cannot capture it, so we offer two ``approximate pictures'' and ask you to imagine their average. That is a limitation of our drawings, not molecular behavior.

\section{How Scientists Know: From Cubic Atoms to Quantum Calculations}

\begin{zhengju}
Covalent bonding was born in an age of ``ledgers without theory.'' In 1916, Lewis proposed that atoms could reach outer octets by sharing electron pairs. He did not yet understand atomic interiors---electron clouds and orbitals lay a decade ahead. His model was even a static ``cubic atom,'' with eight electrons perched at a cube's corners; two cubes sharing an edge shared a pair. Today this looks almost comical, yet it first brought many facts under one explanation: why hydrogen, oxygen, and nitrogen form diatomic molecules, why carbon forms long chains, and why molecular compositions follow simple whole-number ratios. A good model need not start out correct; first it must be \textbf{useful}.

But Lewis could not explain why sharing lowered energy. In his picture, two negative electrons would merely repel each other, so how could sharing them bind atoms? Critics had every reason to object. Another popular theory attributed everything to electrostatics (ionic theory). It explained salt successfully but fell silent before two identical hydrogen atoms: with identical charges on both sides, where was the electrostatic attraction?

The answer awaited quantum mechanics. In 1927, Heitler and London applied the two-year-old Schr\"odinger equation to hydrogen. They wrote the two atoms' electron wavefunctions and calculated total energy against internuclear separation. A valley genuinely emerged, with its position (about \ud{74}{10^{-12}m}) and depth (about \ud{432}{kJ\,mol^{-1}}) agreeing remarkably with measurements. More importantly, the main energy reduction arose from an ``exchange term'' absent from classical physics. As clouds overlap, electrons cannot be distinguished by ``which nucleus they belong to''; states with their positions exchanged must be included in the superposition. This purely quantum effect lowers the energy. A covalent bond is neither a stick nor simple electrostatics, but a clause written into molecules by quantum mechanics. For the first time, chemical bonding became something calculable from equations rather than an empirical mnemonic.

Pauling took the next leg of the relay. He translated quantum mechanics into practical chemical tools: hybridization, resonance, and the electronegativity scale proposed in 1932---our ruler for predicting electron-cloud shifts. His book The Nature of the Chemical Bond influenced generations. But the story also offers a warning: late in life, Pauling became overconfident about vitamin C and other issues, going beyond the evidence. Even the greatest toolbox must have every tool individually tested against evidence.
\end{zhengju}

\section{When These Tools Fail}

Lewis structures are among the most powerful tools a middle-school student can master, but they are \textbf{ledgers, not photographs}. Sooner or later you will encounter their boundaries:

\begin{itemize}[itemsep=2pt]
\item \textbf{Odd-electron molecules cannot have perfect structures}. Nitric oxide, \chem{NO}, has 11 outer electrons, leaving one ``single'' electron however you draw it. This is not your mistake: the molecule really has an unpaired electron. Such molecules (radicals) are particularly reactive, and your body constantly produces and removes them (Chapter 53 returns to them).
\item \textbf{Some atoms prefer fewer than eight}. In boron trifluoride, \chem{BF_3}, boron has only 6 surrounding electrons yet remains stable. Electron deficiency does not mean instability: energy, not a mnemonic, decides.
\item \textbf{Some atoms tolerate more than eight}. Sulfur has 12 surrounding electrons in sulfur hexafluoride, \chem{SF_6}. The octet rule becomes increasingly unreliable from the third period downward.
\item \textbf{Resonance molecules cannot be drawn as they really are}. Ozone, nitrate, and benzene rings have real structures that are ``averages'' of multiple Lewis drawings; any single drawing lies. (Benzene's story is in Chapter 41.)
\item \textbf{Metals ignore this system entirely}. Copper and iron do not distribute electrons pair by pair, but spread them through a ``public sea''---Chapter 21's metallic bonding.
\end{itemize}

The most important boundary also completes this chapter's title: shared electrons \textbf{are not two beads fixed between atoms}. A Lewis line or a pair of dots is only an accounting symbol. The reality is Chapter 9's electron cloud, with electrons spread probabilistically around both nuclei and no definite answer to ``where is this electron right now?'' Covalent bonding is the lowering of total energy after electron clouds overlap. Beads, sticks, and lines are all convenient human representations.

\begin{wuqu}
``A covalent bond means each atom contributes one electron and shares fairly, with neither losing out.'' Two errors bundled together. First, sharing need not be equal: different electronegativities shift the cloud toward the stronger partner. Oxygen pulls shared electrons closer in water, making its end negative and the hydrogen ends positive. A charged comb bending water directly demonstrates that ``sharing is unfair'' (see the experiment). Second, ``one electron each'' is no law: in a coordinate bond, one partner can supply the whole pair. Carbon monoxide can be viewed as oxygen ``contributing'' an extra pair, explaining its unusual formal charge. Likewise, the fourth N--H bond in ammonium, \chem{NH_4^{+}}, uses a pair supplied entirely by nitrogen, yet after forming is completely indistinguishable from the other three. Electrons obey energy, never ``fairness.''
\end{wuqu}

\section{Back to the Blue Cylinder}

Now we can reinterpret everything in the opening.

Oxygen in the cylinder: neither oxygen atom can steal the other's electrons, so they share two pairs, making an \chem{O=O} double bond with energy about \ud{498}{kJ\,mol^{-1}}. Strong enough for compressed storage over years, yet weak enough for a precise molecular machine in your cells to dismantle it gradually and power your entire energy supply (Chapter 53).

The glass of water: within each molecule, polar covalent bonds join oxygen and hydrogen, with the cloud shifted toward oxygen. One end is negative and the other positive, so a charged comb attracts it. Boiling does not dismantle water molecules (that takes thousands of degrees); it only supplies enough energy to escape their loose mutual attractions (Chapters 22 and 23). Boiled water is therefore still water.

The plastic pen barrel and your hair: these are long-chain molecules of carbon, hydrogen, oxygen, nitrogen, and other elements joined covalently. A chain may contain thousands or tens of thousands of covalent bonds. When pulled, chains first slide and straighten, so plastic deforms before breaking. Chapter 43 explores long-chain molecules. A rubber band stretching fivefold and springing back likewise relies on a carbon-chain framework built from covalent bonds.

From cylinders to strands of hair, the same rule supports it all: electron clouds overlap, total energy falls, and atoms join.

\section{After Joining, They Must Arrange Themselves}

This chapter answered how atoms share electrons, but another question immediately arises: shared pairs and lone pairs are both negative, so how do they arrange themselves within a molecule?

Think of water: oxygen has two O--H bonds and two lone pairs. These four clouds repel one another and will not lie in a straight line, so water is bent. Carbon dioxide's two double bonds move as far apart as possible, making \chem{CO_2} straight. Methane's four bonds spread toward a regular tetrahedron's corners. Molecular shape is not random; it follows from electron-pair repulsion. Shape then directly determines what molecules can do. Water's bend exposes its positive and negative ends, while straight \chem{CO_2} makes the two polar bonds cancel. Smell, enzyme catalysis, and DNA's double helix all depend on precise molecular fits. That is next chapter's subject: \textbf{a molecule's shape determines what it can do}.

\begin{tiaozhan}
Why do covalent bonds always share ``pairs,'' not groups of three or five electrons? The answer lies in a quantum property called spin. An electron resembles a tiny magnetic needle with only two possible orientations (up or down), and the Pauli exclusion principle permits at most two electrons in one ``room'' (orbital), with opposite spins. Heitler and London's calculation shows that exchange lowers energy into a valley only for opposite spins; identical spins instead raise energy and push atoms apart. Bonding electrons therefore naturally come in pairs. Going further, another theory beyond Pauling's---molecular orbital theory, developed by Mulliken and others---does not assign electrons to particular bonds at all, but places them throughout the whole molecule. It easily explains facts Lewis theory cannot, such as oxygen molecules actually having two unpaired electrons (liquid oxygen is attracted by magnets!), while Lewis's \chem{O=O} pairs everything. Which theory is right? Both: both are approximations suited to different situations, as scientific models normally are. Skipping this section does not affect the main thread.
\end{tiaozhan}

\section*{Try It}

\begin{sikaoti}
\item Draw ammonia's Lewis structure, \chem{NH_3} (nitrogen brings 5 outer electrons and each hydrogen 1). Count nitrogen's bonding pairs and lone pairs. What is the total electron budget, and have you used it all?
\item Explain energetically why ``breaking covalent bonds absorbs energy,'' not releases it. Draw a two-atom potential-energy curve, with energy vertically and internuclear separation horizontally, and mark bond length and bond energy.
\item Look up electronegativities and rank these bonds by increasing electron-cloud shift: C--H, O--H, H--H, Na--Cl, H--Cl. Which pair lies beyond the covalent range and is more ionic?
\item Someone says, ``A double bond is two single bonds, so it should be twice as strong.'' Test this against the table's \chem{C-C} (347) and \chem{C=C} (614) values. If it fails, suggest why the second bond is ``weaker than the first.'' (Hint: two electron pairs between the same atoms repel each other.)
\item Nitrate, \chem{NO_3^{-}}, has three equivalent resonance structures, with the double bond on each oxygen in turn. Predict whether the three N--O bond lengths are equal, and how much negative charge each oxygen shares.
\item Strong magnets attract liquid oxygen, indicating unpaired electrons, yet the Lewis structure \chem{O=O} pairs every electron. What does this contradiction reveal about Lewis structures as a tool? Should we discard them or redefine their role?
\end{sikaoti}
